\documentclass[11pt,a4paper]{article}

\usepackage[margin=2.5cm]{geometry}
\usepackage{amsmath,amssymb}
\usepackage{bm}
\usepackage{graphicx}
\usepackage{hyperref}
\usepackage[numbers]{natbib}

\newcommand{\e}{{\mathrm{e}}}
\renewcommand{\i}{{\mathrm{i}}}
\newcommand{\d}{{\mathrm{d}}}

\title{Learning to Remember: Error-Driven Attractor Formation\\
in Thermodynamic Dense Associative Memories}
\author{}
\date{}

\begin{document}
\maketitle

%===========================================================
\begin{abstract}
Continual learning systems must acquire new knowledge over time without erasing previously stored information.
A key missing piece in many approaches is an explicit mechanism for \emph{learning when to remember}---i.e., when the system should form a new memory rather than adapt existing representations.
We formulate this problem in Hopfield-type dense associative memories, where memories correspond to attractors of an energy (or free energy) landscape.
Our central hypothesis is that large prediction errors signal the absence of a suitable attractor for the current input and should therefore trigger \emph{attractor formation}.
Using thermodynamic dense associative memories (Thermo-DAMs), we provide an energy-based framework in which memory formation corresponds to controlled deformations of the free energy landscape, enabling principled reasoning about stability, basin geometry, and interference in continual learning.
\end{abstract}

%===========================================================
\section{Introduction}

Continual learning addresses learning from a non-stationary stream of data while preserving previously acquired knowledge.
A central difficulty is catastrophic forgetting: incorporating new information can destroy older knowledge.
Most existing solutions mitigate forgetting through parameter regularization, architectural isolation, or replay of past data.
While effective in certain regimes, these approaches often treat memory as an auxiliary resource and do not explicitly address a fundamental decision problem: \emph{when should a system form a new memory?}

We argue that continual learning is fundamentally a problem of \emph{learning to remember}.
In particular, a learning system operating online must decide \emph{when} to create a new memory trace (or attractor) and when to rely on existing ones.
This decision is nontrivial: indiscriminate memorization leads to interference and uncontrolled growth, while overly conservative updates prevent adaptation to novel structure.

Hopfield-type associative memories provide a natural language for this question.
In dense associative memories (DAMs), stored patterns correspond to attractors of an energy landscape, and retrieval is realized by dynamical convergence.
From this viewpoint, ``remembering'' is not an explicit data storage operation but the existence of a stable basin of attraction that supports robust retrieval.
Consequently, learning to remember can be cast as the problem of \emph{when and how to create new attractors}.

In this work, we propose a principled trigger for attractor formation: \emph{prediction error}.
Intuitively, if an input can be explained by convergence to an existing attractor, prediction error remains small and no new memory is required.
If prediction error becomes large, this suggests that no suitable attractor exists---the system fails to retrieve an internal state that supports correct prediction---and a new attractor should be formed.
We formalize this idea within thermodynamic dense associative memories (Thermo-DAMs), where attractor stability and interference can be analyzed through the free energy landscape and its entropic geometry.

%===========================================================
\section{Related Work}

\paragraph{Continual learning.}
Continual learning has been studied through regularization-based approaches (constraining parameter drift), replay mechanisms (storing and revisiting past samples), and architectural expansion (allocating task-specific subnetworks).
These methods primarily operate at the parameter or data level.
In contrast, our focus is the \emph{memory formation decision}: when a novel event should be committed as a new stable memory state.

\paragraph{Hopfield networks and dense associative memories.}
Hopfield networks and modern dense variants have been analyzed in terms of capacity, retrieval dynamics, and connections to energy-based models and attention.
Most work assumes a fixed set of stored patterns and studies retrieval.
We instead consider the continual regime, where the memory set itself evolves, and ask when new attractors should be created.

\paragraph{Prediction error and surprise signals.}
Prediction error is a central concept in learning theories across neuroscience and machine learning.
In machine learning it is typically used to drive incremental parameter updates.
Here we use it as a \emph{structural signal}: a criterion for creating new attractors in an associative memory.

%===========================================================
\section{Problem Formulation: Learning When to Remember}

We consider a system that processes a stream of inputs $\{\bm{x}_t\}$ and produces predictions for a task of interest (e.g., classification, next-step prediction, or reconstruction).
At each time $t$, the system incurs a prediction error $e_t$ (task-dependent, e.g., negative log-likelihood or squared error).

The system is equipped with a Hopfield-type dense associative memory with stored patterns $\{\bm{\xi}^\mu\}_{\mu=1}^p$.
Memory retrieval corresponds to dynamical convergence to an attractor (a local minimum of an energy or free energy function) given an initial state or cue.

\paragraph{Key question.}
The central problem is to learn a rule for \emph{memory formation}:
\begin{quote}
\emph{When should the system modify its attractor landscape by creating a new stable memory state (attractor), rather than relying on existing attractors or incremental adaptation?}
\end{quote}

\paragraph{Error-driven attractor formation.}
Our guiding hypothesis is that prediction error provides a principled trigger:
small errors indicate that the current input is well explained by existing attractors;
large errors indicate the absence of a suitable attractor and should trigger \emph{attractor formation}.
Thermo-DAMs provide the appropriate framework to analyze this process in terms of the free energy landscape, basin geometry, and entropic costs of introducing additional attractors.

%===========================================================
\section{Thermodynamic Dense Associative Memory Model}
\label{sec:thermo_dam}

We now introduce the thermodynamic dense associative memory (Thermo-DAM) framework underlying our formulation.
This model provides an energy-based description of stored patterns, retrieval dynamics, and basin geometry.
Crucially, it also enables reasoning about how the landscape changes when new memories are formed.

% -------------------- YOUR TEXT STARTS HERE (UNCHANGED) --------------------
% Insert your full Thermo-DAM derivations and equations here as-is.
% -------------------- YOUR TEXT ENDS HERE --------------------

%===========================================================
\section{Error-Driven Attractor Formation in Thermo-DAMs}
\label{sec:formation}

In Hopfield-type associative memories, successful retrieval corresponds to convergence toward an existing attractor.
If an input lies within the basin of attraction of a stored memory, the system reaches a stable state that supports correct prediction, and the prediction error remains small.

When prediction error is large, the system fails to find a stable internal state compatible with the input and the task objective.
In the language of energy-based memories, this suggests the absence of an appropriate attractor in the relevant region of state space.
We propose to interpret such events as \emph{signals to remember}: the system should form a new attractor that encodes the novel structure revealed by the error and its subsequent correction.

Within the thermodynamic formulation, memory formation corresponds to a controlled deformation of the free energy landscape.
Introducing a new attractor increases the set of retrievable states but incurs an entropic cost by reducing the available volume of microstates compatible with a given set of alignments.
This naturally yields a trade-off between adaptability (forming new memories) and stability (preserving existing basins), providing a principled viewpoint on interference and forgetting.

Importantly, this perspective does not require proposing a new architecture.
Thermo-DAMs serve as a theoretical and computational substrate for formulating continual learning as \emph{learning when to remember} via error-driven attractor formation.

%===========================================================
\section{Experimental Protocol}
\label{sec:protocol}

We outline experimental protocols aimed at diagnosing whether error-driven attractor formation improves continual learning.
The emphasis is on testing the causal chain:
\[
\text{large prediction error} \;\Rightarrow\; \text{attractor formation} \;\Rightarrow\; \text{improved future retrieval and retention}.
\]
Candidate settings include class-incremental learning, noisy pattern completion, and task-switching.
In each case, the key diagnostic is whether error-triggered memory formation leads to stable attractors with basins that support robust retrieval under future perturbations.

%===========================================================
\section{Discussion}

Our formulation reframes continual learning as the problem of shaping an attractor landscape online.
In this view, forgetting is not merely parameter drift but the loss of basin structure supporting retrieval of earlier memories.
Prediction error becomes a structural signal for deciding \emph{when} a new memory should be formed.

The thermodynamic viewpoint adds an essential layer: it provides explicit control variables (temperature, entropic penalties) and yields interpretable trade-offs between memory growth, stability, and robustness.
This perspective suggests new evaluation criteria for continual learning beyond accuracy curves, such as basin volumes, attractor stability, and free energy barriers.

%===========================================================
\section{Conclusion}

We proposed a Hopfield-centric formulation of continual learning as \emph{learning to remember}: deciding when to form new memories.
We argued that large prediction errors signal the absence of a suitable attractor and should trigger error-driven attractor formation.
Thermodynamic dense associative memories provide a principled framework for analyzing and implementing this mechanism through free energy landscape shaping.
This viewpoint opens new directions for studying continual learning as a dynamical memory formation problem grounded in attractor geometry.

\bibliographystyle{plainnat}
% \bibliography{refs}

\end{document}
